

% ---- From Santi (Removed by Carlos)
% MadGraph5\_aMC@NLO~\cite{Alwall:2014hca,Frederix:2018nkq} (MG5aMC) is a state-of-the-art framework widely used in high-energy physics for the automatic generation of matrix elements and event simulation at leading (LO) and next-to-leading order (NLO) in perturbative quantum field theory. It provides a modular, extensible, and community-supported platform that is actively being adapted to meet the computational challenges posed by the HL-LHC and future colliders.

% Although experimental uncertainties today are larger than those of simulations, at the High Luminosity LHC experimental precision is expected to be above the theoretical one for events generated below NLO precisions. As forecasted hardware resources will not meet CPU requirements for these simulation needs, speeding up NLO event generation is a necessity.

% The integration of GPU acceleration has demonstrated substantial performance improvements. For complex processes such as $gg\rightarrow t\bar{t}ggg$, the achieved speed-up on the ME calculation utilization of NVIDIA V100 GPUs has resulted in ME computation ranging from 100x to 1000x compared to single-threaded Fortran implementations. This acceleration translates to an overall workflow speed-up of approximately 60x to 100x, considering the entire event generation pipeline. Additionally, the use of vector instructions on CPUs has led to ME calculation speed-ups of up to 16x. When combined with multi-core processing, the overall event generation workflow experiences a performance enhancement of approximately 6x to 10x.

% NLO computations introduce additional complexities, including real emission processes and loop corrections, which inherently involve more intricate algorithms and data dependencies compared to LO calculations. To address these challenges, the MG5aMC development team is implementing a multi-event processing framework that facilitates vectorization and parallel execution. This approach aims to exploit the capabilities of modern hardware architectures, such as GPUs, to accelerate the evaluation of NLO amplitudes. Preliminary efforts have demonstrated the feasibility of this strategy, with ongoing work focusing on optimizing the computational kernels and integrating them seamlessly into the existing MG5aMC infrastructure~\cite{Wettersten:2025hrb}. 

% On this work we will focus on the generation of the simple $gg\rightarrow t\bar{t}g$ process. It is true that in this process the non-ME part is dominating the CPU computation, but the ME-part is simple enough that allows for a simple case-study that we show in this manuscript. 


Among the most widely used frameworks for MC event generation in high-energy physics is the \amcnlo~(\mgfive) software. The team has already made significant strides in acceleration methods, particularly focusing on GPU and vector CPU support. The release of the first production version of the CUDACPP plugin (v1.00.00) ~\cite{Valassi:2025xfn} enables the acceleration of LO processes using GPUs and vector CPUs and involves extensive re-engineering to support data-parallel execution, facilitating significant speed-ups in ME computations, which traditionally represent the primary computational bottleneck in event generation workflows. 

The integration of GPU acceleration has demonstrated substantial performance improvements. For complex processes, the achieved speed-up on the ME calculation utilization of NVIDIA V100 GPUs has resulted in ME computation ranging from 100x to 1000x improvements compared to single-threaded Fortran implementations. This acceleration translates to an overall workflow speed-up of approximately 60x to 100x, considering the entire event generation pipeline. Additionally, the use of vector instructions on CPUs has led to ME calculation speed-ups of up to 16x. When combined with multi-core processing, the overall event generation workflow experiences a performance enhancement of approximately 6x to 10x.

The \mgfive~framework provides several features that make it particularly suitable for use as a benchmark when testing the performance of matrix element integration on hardware-based accelerators such as FPGAs. One such feature is the high degree of flexibility available to users to configure physics processes, model parameters, and integration settings. Once these parameters are fixed, the \mgfive~solution provides an automatic way of performing the matrix element integration by taking care of producing the necessary code to run the actual calculations. As a result, evaluating the limits of resource usage, throughput, numerical precision, and latency of an FPGA-based solution can be practically achieved by porting the matrix element code generated by \mgfive~to a hardware description language (HDL) of choice. A key advantage of this approach is that it removes the need for in-depth knowledge of the specific physics calculations being performed, since \mgfive~already provides validated and optimized implementations of the required matrix elements. This allows the focus to be placed directly on the hardware performance and architectural aspects of the accelerator, rather than on the underlying theoretical formulation of the problem.

As this work represents a first attempt at implementing complex physical calculations on a hardware accelerator, the focus is placed on exploring the limitations of FPGAs in one of the simplest scenarios that can be realized in hardware: the production of a \ttbar~final state through the interaction of two gluons from the incoming colliding protons, namely \(\ggtt\). The accuracy of the calculation is restricted to the LO perturbative expansion, since the complexity of NLO calculations can be challenging even for GPU-based solutions. To increase the complexity of the problem and provide a more demanding benchmark for the FPGA, we nevertheless include the first real-emission corrections, which constitute a subset of the full set of contributions required to achieve true NLO accuracy. To this end, we instruct \mgfive~to generate the code for the \ggttj~process, where the $\mathrm{j}$ denotes an additional emission from either the initial-state (\gluglu) or final-state (\ttbar) particles.

Within \mgfive, the complete integration of the matrix element is factorized into different steps. The most demanding step for processes such as \ggttj correspondes to the computation of the color algebra [...] % need to fill the gap here.



